\section{Preliminaries and imported ingredients}\label{sec:preliminaries}

\subsection{Notation}

All logarithms are natural unless a base is displayed. Constants denoted by $c,C,C'$ are positive and may change from line to line. A constant carrying a subscript may depend on that subscript but not on $n$.

Fix a parameter $\sigma\in(0,1)$. Throughout the structural argument in \cref{sec:profiles,sec:kernels,sec:structural,sec:normals,sec:intermediate} and the lower-band hermitization argument, the integer degree satisfies
\begin{equation}\label{eq:lower-band}
  C_*\log n\le d\le(\log n)^{1+\sigma}.
\end{equation}
All constants in those arguments are uniform in $d$ on this interval; the threshold from which the estimates hold may depend on $\sigma$. The superlogarithmic bridge in \cref{subsec:superlog-bridge} is stated separately and does not assume~\eqref{eq:lower-band}. Rounding is suppressed only inside asymptotic relations such as $m=n/d^3+O(1)$.

\medskip
\noindent\textbf{Constant convention.}
Let $a_3^{\src},\cK,\cP,c_{4.10}>0$ denote fixed admissible constants from the cited LLTTY statements. Choose once and for all
\[
  0<a_3\le a_3^{\src}
\]
sufficiently small for the entropy and row-load inequalities in \cref{subsec:shifted-T1T2,subsec:almost-constant}. This is permitted because the source construction takes $a_3$ to be a sufficiently small absolute constant; throughout the manuscript the LLTTY classes and the scale $n_3=\floor{a_3n}$ are defined using this single final choice. Next choose the absolute normal parameter $a\le c_{4.10}$, and then enlarge the absolute constant $C_*$ so that the support-diameter estimate, the imported smallest-singular-value and restricted-norm estimates, and the rectangular witness and net estimates all hold. After these absolute choices, fix $\sigma\in(0,1)$ and a shift bound $Z_0<\infty$. Only the threshold
\[
  n\ge n_0(\sigma,Z_0)
\]
may depend on these two parameters. No later argument further decreases $a_3$ or enlarges a source constant.

For a vector $x\in\C^n$, let
\[
  x_1^*\ge x_2^*\ge\cdots\ge x_n^*
\]
be the nonincreasing rearrangement of $(|x_j|)_{j\le n}$. For a matrix $H$ with at least as many columns as rows, its singular values are denoted
\[
  s_1(H)\ge\cdots\ge s_{\min\{\#\text{rows},\#\text{columns}\}}(H)\ge0.
\]
For a $k\times n$ matrix $H$, define its row volume by
\[
  \vol(H):=\sqrt{\det(HH^*)}=\prod_{j=1}^k s_j(H).
\]
We use
\[
  \logp t:=\max\{\log t,0\},
  \qquad
  \logm t:=\max\{-\log t,0\}.
\]

For probability measures $\mu,\nu$ on $\C$, write
\[
  d_{\BL}(\mu,\nu)
  :=\sup\set{
    \abs{\int f\,d\mu-\int f\,d\nu}:
    \norm{f}_\infty\le1,
    \ \Lip(f)\le1
  }.
\]
For random probability measures, $\mu_n\Rightarrow\mu$ in probability means $d_{\BL}(\mu_n,\mu)\to0$ in probability. The same convention is used for differences of two random empirical spectral measures.

Our Hermitian inner product is
\[
  \inner{u}{v}=\sum_{j=1}^n u_j\overline{v_j},
\]
linear in the first variable. Matrix rows act bilinearly:
\[
  r\cdot x=\sum_{j=1}^n r_jx_j.
\]
Thus $r\cdot x=\inner{x}{\overline r}$. This convention is used throughout the right-kernel, Gaussian-normal, and leave-one-row-out sections.

A row slice is a random vector $X$ uniformly distributed on $\Omega_{n,d}$. We write $p=d/n$. A Bernoulli comparator row is denoted by
\[
  \boldsymbol\xi=(\xi_1,\ldots,\xi_n),
  \qquad
  \xi_j\stackrel{\mathrm{iid}}\sim\Ber(p).
\]
We also write
\[
  \C^\times:=\C\setminus\{0\},
  \qquad
  \C^\times\mathcal X:=\set{\lambda x:\lambda\in\C^\times,\ x\in\mathcal X}
\]
for the nonzero scalar cone generated by a set $\mathcal X\subset\C^n$.

\subsection{\texorpdfstring{Discretization and the LLTTY $\ell$-decomposition}{Discretization and the LLTTY ell-decomposition}}

We recall only the portions of the discretization from~\cite[Sections~5--7]{LLTTY19} that enter our proof.

\begin{definition}[$k$-vector]\label{def:k-vector}
For an integer $k\ge1$, a vector $y\in\C^n$ is a $k$-vector if $ky_j\in\Z+i\Z$ for every $j$. The canonical $k$-approximation of $x\in\C^n$ is
\[
  y_j=\frac1k\bigl(\floor{k\Re x_j}+i\floor{k\Im x_j}\bigr).
\]
\end{definition}

Given a $k$-vector $y$, its coordinates are partitioned into level sets and then into regular and spread $\ell$-parts by the source construction. At order $j\ge0$, let $I_j$ be the coordinates not assigned at an earlier order and let $\Lambda_j=\{y_i:i\in I_j\}$. For $\lambda\in\Lambda_j$, the level set $L(j,\lambda)$ is the entire remaining $\lambda$-level when that level has cardinality less than $2^{j+1}$, and otherwise is its leftmost subset of cardinality $2^j$. Thus every nonempty $L(j,\lambda)$ has cardinality between $2^{j-1}$ and $2^{j+1}$.

The spread values at order $j$ are chosen as follows. If
\[
  \diam(\Lambda_j)<d/k,
\]
then the spread family is empty and every level set of order $j$ is assigned to the regular part. Otherwise the canonical greedy procedure of~\cite[Section~5.2]{LLTTY19} selects a maximal pairwise $d/k$-separated family $\Lambda_j^S\subset\Lambda_j$, starting from a separated pair. In particular, $\Lambda_j^S$ then contains at least two values. The spread and regular parts of order $j$ are
\[
  \mathcal L_j^S=\bigcup_{\lambda\in\Lambda_j^S}L(j,\lambda),
  \qquad
  \mathcal L_j^R=\bigcup_{\lambda\in\Lambda_j\setminus\Lambda_j^S}L(j,\lambda).
\]
Consequently, a singleton family of represented values is always regular, every nonempty spread part has height at least two, and the distinct values represented in a spread part are pairwise $d/k$-separated. Every value represented in a regular part is either in a set of diameter less than $d/k$ (when $\mathcal L_j^S=\varnothing$) or lies at distance less than $d/k$ from the greedy spread family.

Empty parts are omitted and the nonempty spread and regular parts are re-enumerated as
\[
  \mathcal L^{(1)},\ldots,\mathcal L^{(r)},
  \qquad r\le3\log n.
\]
The height $h_q=h(\mathcal L^{(q)})$ is the number of level sets comprising $\mathcal L^{(q)}$. If $\mathcal L^{(q)}$ has order $j$, then
\begin{equation}\label{eq:part-size-height}
  2^{j-1}h_q\le s_q:=|\mathcal L^{(q)}|\le2^{j+1}h_q.
\end{equation}
We use the source's lexicographic tie-breaking and ordering of the nonempty parts, which make the decomposition canonical. After permuting coordinates, the types, orders, sizes, and heights of the parts agree up to relabeling. Only these invariants are used below.

For $p_q=s_q/n$ and a profile count $Q_q\in\{0,\ldots,d\}$, define the actual weights
\begin{equation}\label{eq:actual-weights}
  w_q(Q)=
  \begin{cases}
    h_qQ_q/d,&\mathcal L^{(q)}\text{ regular},\\
    h_q\sqrt{Q_q},&\mathcal L^{(q)}\text{ spread},
  \end{cases}
\end{equation}
and the truncated weights
\begin{equation}\label{eq:truncated-weights}
  \wt w_q=
  \begin{cases}
    h_qp_q,&s_q\ge nd^{-1/3},\ \mathcal L^{(q)}\text{ regular},\\
    h_q/d,&s_q<nd^{-1/3},\ \mathcal L^{(q)}\text{ regular},\\
    h_q\sqrt{dp_q},&s_q\ge nd^{-1/3},\ \mathcal L^{(q)}\text{ spread},\\
    h_q,&s_q<nd^{-1/3},\ \mathcal L^{(q)}\text{ spread}.
  \end{cases}
\end{equation}
In particular,
\begin{equation}\label{eq:weight-range}
  \frac{h_q}{d}\le\wt w_q\le n\sqrt d.
\end{equation}
The dyadic $w$-set of order $b$ is the union of those $\mathcal L^{(q)}$ for which
\[
  2^b\le\wt w_q<2^{b+1}.
\]
The conditional one-row estimator is
\begin{equation}\label{eq:SB-actual}
  \operatorname{SB}(Q):=\min\left(1,\min_{q:Q_q>0}w_q(Q)^{-1}\right),
\end{equation}
and its truncated analogue is
\begin{equation}\label{eq:SB-truncated}
  \operatorname{SB}(\wt Q):=\min\left(1,\min_{q:Q_q>0}\wt w_q^{-1}\right).
\end{equation}

\subsection{Author-derived deterministic consequences of the LLTTY decomposition}

The following result is extracted from~\cite[Lemmas~5.1, 5.4, 5.5, 7.2, 7.7--7.9 and Theorem~5.6]{LLTTY19}. Its $K_u$ and clustered-$P_v$ conclusions include intermediate estimates from the proofs of the cited source lemmas, not merely verbatim source statements. Appendix~\ref{app:LLTTY-extraction} gives a lemma-by-lemma derivation, with exact source-version, page, lemma, and equation locators. The source works in the degree range $d\le\exp(c\sqrt{\log n})$ for an absolute $c>0$. Under~\eqref{eq:lower-band}, every fixed polylogarithmic upper bound is eventually smaller than this quantity, so the global source restriction is automatic.

In the lower-band argument, set
\[
  m:=\floor{n/d^3}-1
\]
and call a codimension admissible if
\[
  L\in\{m,m+1\}
  =\set{\floor{n/d^3}-1,\floor{n/d^3}}.
\]
For all sufficiently large $n$, one has $m\ge1$, and every admissible codimension satisfies the source restriction
\[
  1\le L\le\frac n{d^3}.
\]

For each admissible codimension $L$, let $S_L$ denote the exact normalized gradual class of~\cite[Section~4]{LLTTY19}, formed with the codimension-dependent steep classes and the final constant $a_3$ fixed above. This is the source class used in \cref{sec:structural,sec:normals,sec:intermediate}; \cref{subsec:classes} reproduces its defining inequalities in full. Thus no statement below depends on a definition deferred to a later section.

For $x\in\C^n$, a scale $u>0$, and $\rho>0$, a set $J\subset[n]$ is a $\rho$-cluster of $x$ at scale $u$ if there is $\lambda\in\C$ such that
\[
  |x_j-\lambda|\le\rho u\qquad(j\in J).
\]
When $u=x^*_{n_3}$, we simply say that $x$ has a $\rho$-cluster of cardinality $|J|$.

For an integer $k\ge1$, let
\[
  \mathcal A_k(L):=\set{y:y\text{ is the canonical $k$-approximation of some }x\in S_L}.
\]
Fix the lexicographic order on $\C$, let $y^\sharp$ denote the coordinates of $y$ rearranged in nonincreasing lexicographic order, and, for $y,\wt y\in\mathcal A_k(L)$, set
\[
  y\sim_k\wt y\quad\Longleftrightarrow\quad y^\sharp=\wt y^\sharp.
\]
Thus an equivalence class consists precisely of normalized canonical approximants having the same coordinate multiset; this relation is not asserted on the unrestricted, infinite set of all $k$-vectors. For $\mathcal X\subset S_L$, write
\[
  \mathcal A_k(\mathcal X)
  :=\left\{y:\begin{array}{l}
    y\text{ is the canonical $k$-approximation}\\[-2pt]
    \text{of some }x\in\mathcal X
  \end{array}\right\}
  \subset\mathcal A_k(L).
\]
and use the restriction of $\sim_k$ to this set.

\begin{proposition}[Author-derived deterministic consequences of the LLTTY decomposition]\label{prop:LLTTY-package}
There are absolute constants $a_3,\cK,\cP,c_0,c_1>0$ such that, for every admissible codimension $L$ and the exact normalized gradual class $S_L$, there are families $K_u(L),P_u(L)\subset S_L$ with the following properties.
\begin{enumerate}[label=\textup{(\roman*)}]
\item For every integer $v\ge5$,
\[
  S_L=\bigcup_{u=4}^{v}K_u(L)\cup P_v(L).
\]

\item Suppose additionally that $d\le n^{1/3}$ and
\[
  1\le k\le\frac{\sqrt n}{d^{3/2}}.
\]
Then $\mathcal A_k(L)$ has at most $e^n$ equivalence classes under $\sim_k$. The same bound holds for the nonempty classes obtained by restricting $\sim_k$ to any $\mathcal A_k(\mathcal X)\subset\mathcal A_k(L)$. If $\mathcal C$ is such a class with $\ell$-part sizes $s_q$ and heights $h_q$, then
\begin{equation}\label{eq:class-cardinality}
  |\mathcal C|\le\frac{n!}{\prod_qs_q!}\prod_q h_q^{s_q}.
\end{equation}

\item If $I(b)=\{q:2^b\le\wt w_q<2^{b+1}\}$ and $t_b=\sum_{q\in I(b)}s_q$, then
\begin{equation}\label{eq:factorial-compression}
  \prod_{q\in I(b)}s_q!\ge t_b!e^{-Ct_b}.
\end{equation}
Moreover, for $J(b)\subset I(b)$ consisting of parts of size $<d^{-\kappa}n$, where
\[
  \kappa=\frac{\cK a_3}{16}<\frac13,
\]
one has
\begin{equation}\label{eq:small-bin-mass}
  \left|\bigcup_{q\in J(b)}\mathcal L^{(q)}\right|\le Cd^{-\kappa}n.
\end{equation}

\item If $y$ is a $d^u$-approximation of a vector in $K_u(L)$, then the total mass of spread $\ell$-parts is at least $\cK n_3$, where $n_3=\floor{a_3n}$. With
\[
  I=\{q:s_q\ge d^{-\kappa}n\},
\]
one of the two alternatives holds:
\begin{align}
  \sum_{q\in I}s_q&\ge n-\frac{\cK n_3}{4},\tag{K1}\label{eq:K1}\\
  \sum_{q\in I^c}s_q&\ge\frac{\cK n_3}{4}.\tag{K2}\label{eq:K2}
\end{align}
In alternative~\eqref{eq:K1},
\begin{equation}\label{eq:K1-product}
  \prod_qh_q^{s_q}\prod_q\wt w_q^{-s_q}\le C^n d^{-c_1n}.
\end{equation}
In alternative~\eqref{eq:K2}, there are integers $b_1\le b_2$ such that the coordinate distribution of the dyadic weight order has mass at least $c_0$ below $b_1$ and at least $c_0$ above $b_2$, and
\begin{equation}\label{eq:K2-separation}
  b_2-b_1\ge c_1d^\kappa.
\end{equation}

\item Fix $v\ge5$ and $\rho\le d^{-v}$. Assume that
\begin{equation}\label{eq:Pv-delta}
  0<\delta\le36\cP a_3
\end{equation}
and
\begin{equation}\label{eq:Pv-separation-hyp}
  \cP a_3(v-4)
  \ge2\log_2d+2\log_2\!\left(\frac{72\sqrt d}{\delta}\right)
     +2-\log_2(\cP a_3).
\end{equation}
If $x\in P_v(L)$ has a $\rho$-cluster of at least $\delta n$ coordinates and $y$ is its $d^v$-approximation, then the dyadic weight-order distribution has mass at least $\delta/36$ below some $b_1$, mass at least $c_0$ above some $b_2$, and
\begin{equation}\label{eq:Pv-separation}
  b_2-b_1\ge\frac{\cP a_3}{2}(v-4).
\end{equation}
\end{enumerate}
\end{proposition}

The source status of the five parts is summarized below. Page numbers refer to arXiv:1801.05575v4.
\begin{center}
\small
\begin{tabular}{p{0.07\textwidth}p{0.36\textwidth}p{0.47\textwidth}}
\toprule
Part & Status & Exact locator\\
\midrule
(i) & verbatim decomposition after notation change & p.~38, Theorem~5.6\\
(ii) & direct restatement and reindexing & pp.~34--35, Lemma~5.1 and (14)--(16); p.~37, Lemmas~5.4--5.5\\
(iii) & stated factorial bound plus proof-level cutoff consequence & p.~55, Lemma~7.2; p.~61, Case~2 in the proof of Lemma~7.7\\
(iv) & author-derived packaging of the two proof branches & pp.~60--61, proof of Lemma~7.7, equations (32)--(33), (38), Cases~1--2\\
(v) & author-derived combination of two named lemmas & p.~62, Lemma~7.8; pp.~62--63, Lemma~7.9 and its proof\\
\bottomrule
\end{tabular}
\end{center}

\begin{remark}\label{rem:source-map-package}
Appendix~\ref{app:LLTTY-extraction} shows in particular that \eqref{eq:K1-product} is the three-family calculation in the first case of~\cite[Lemma~7.7]{LLTTY19}, while the separated weight populations in~\eqref{eq:K2} and part~(v) are intermediate conclusions used in the proofs of~\cite[Lemmas~7.7 and 7.9]{LLTTY19}. The factorial estimate~\eqref{eq:factorial-compression} and the small-bin estimate~\eqref{eq:small-bin-mass} use exactly the source's truncated weights and $w$-sets; no change of binning is made.
\end{remark}

\subsection{Row-slice inputs from Li--Litvak--Yu}

We use the following consequences of~\cite{LLY26a}.

\begin{proposition}[Smallest singular values]\label{prop:LLY-smallest}
There are absolute constants $C,C_1,C_2>0$ such that, whenever $C\log n\le d\le n/2$ and $|\zeta|\le\sqrt d\log\log d$, both the independent row-slice matrix $M_n$ and the i.i.d. Bernoulli$(d/n)$ matrix $B_n$ satisfy
\[
  \Pp\set{s_n(M_n-\zeta I_n)\le n^{-\beta_n}}
  +\Pp\set{s_n(B_n-\zeta I_n)\le n^{-\wt\beta_n}}
  \le\frac{C_1}{\sqrt d},
\]
for exponents
\[
  2<\beta_n,\wt\beta_n\le C_2\log\log n.
\]
\end{proposition}

This combines~\cite[Theorem~1.7, Remark~1.9, Theorem~3.1 and Remark~3.2]{LLY26a}.

\begin{corollary}[Uniform polynomial scale above a superlogarithmic cutoff]\label{cor:uniform-polynomial}
Fix $\eps>0$. There exists a constant $B_\eps<\infty$ such that, whenever
\[
  (\log n)^{1+\eps}\le d\le n/2,
  \qquad
  |\zeta|\le\sqrt d\log\log d,
\]
one has
\begin{equation}\label{eq:uniform-polynomial}
  \Pp\set{s_n(M_n-\zeta I_n)\le n^{-B_\eps}}
  +\Pp\set{s_n(B_n-\zeta I_n)\le n^{-B_\eps}}
  \le\frac C{\sqrt d}.
\end{equation}
\end{corollary}

\begin{proof}
We make the extraction from the source explicit. Write
\[
  \Gamma(n,d):=\frac{2\log(6d)}{\log(d/(C_2\log n))}.
\]
In the sparse branch of~\cite[Remark~1.9]{LLY26a}, the admissible threshold, denoted there by $\alpha(n,d)$, satisfies
\[
  \alpha(n,d)\ge n^{-2\Gamma(n,d)}.
\]
If $d\ge(\log n)^{1+\eps}$, then for all sufficiently large $n$,
\[
  \log\frac d{C_2\log n}
  =\log d-\log\log n-\log C_2
  \ge\frac{\eps}{2(1+\eps)}\log d.
\]
Also $\log(6d)\le2\log d$ for large $d$. Consequently,
\[
  \Gamma(n,d)\le\frac{8(1+\eps)}{\eps}
\]
uniformly throughout the sparse branch. Hence the source threshold there is at least $n^{-B_\eps}$ for any fixed
\[
  B_\eps\ge\frac{16(1+\eps)}{\eps}.
\]
In the complementary dense branch, the same remark gives
\[
  \alpha(n,d)\ge\frac1{2nd^{3/2}}\ge n^{-3}
\]
for all sufficiently large $n$, since $d\le n/2$. Increasing $B_\eps$ if necessary therefore gives a single polynomial threshold throughout $d\le n/2$. Remark~3.2 of the same paper permits exactly the same threshold for the Bernoulli comparator. The probability estimate is the one in \cref{prop:LLY-smallest}.
\end{proof}

\begin{proposition}[Restricted operator norm in the range used below]\label{prop:restricted-norm}
Let $H$ be an $m\times n$ matrix with independent uniform $d$-subset rows. Assume
\[
  \frac n2\le m\le n,
  \qquad
  d\le\frac n2.
\]
For every fixed $A>0$ there is $C_A<\infty$ such that, whenever $d\ge C_A\log n$,
\[
  \norm{H-\E H}\le C_A\sqrt d
\]
with probability at least $1-n^{-A}$.
\end{proposition}

\begin{proof}
Apply~\cite[Theorem~4.1]{LLY26a} with its parameters $\gamma=A+4$ and $\beta=A+1$. The source failure probability is
\[
  nm^{-\gamma}+2n^{-\beta}.
\]
Since $m\ge n/2$, this is at most $n^{-A}$ for all sufficiently large $n$. Write
\[
  N=\max\{9d,2\gamma\log m\}.
\]
Because $m\le n$ and $d\ge C_A\log n$, enlarging $C_A$ if necessary gives $N=9d$. The source estimate
\[
  \norm{H-\E H}
  \le \frac{2d}{\sqrt N}+C_1\beta\sqrt N
\]
therefore yields $\norm{H-\E H}\le C_A\sqrt d$. This is the only range of the rectangular theorem used in the manuscript.
\end{proof}

The right-kernel exclusions needed later are not imported: their complete rectangular proof is given in \cref{sec:kernels}.

\subsection{Concentration and tensorization}

We use the L\'evy concentration function
\[
  \mathcal Q(\xi,t):=\sup_{\lambda\in\C}\Pp\set{|\xi-\lambda|\le t}.
\]
The following standard tensorization statement is the form used in~\cite[Lemma~2.7]{LLTTY19}.

\begin{lemma}[Tensorization]\label{lem:tensorization}
Let $\xi_1,\ldots,\xi_K$ be independent complex random variables. If
\[
  \mathcal Q(\xi_i,\tau t)\le A\tau^2
  \qquad\text{for all }\tau\ge1\text{ and all }i,
\]
then, for an absolute $C$ and every $\tau\ge1$,
\[
  \Pp\set{\sum_{i=1}^K|\xi_i|^2\le C\tau^2t^2K}
  \le(CA\tau^2)^K.
\]
The statement remains valid after deterministic shifts of the $\xi_i$.
\end{lemma}

We also use the negative second-moment identity. If $H$ is a full-row-rank $k\times n$ matrix with rows $R_i$, then
\begin{equation}\label{eq:negative-second-moment}
  \sum_{j=1}^k s_j(H)^{-2}
  =\sum_{i=1}^k\dist\bigl(R_i,\spanv(R_\ell:\ell\ne i)\bigr)^{-2}.
\end{equation}
The inequality needed later remains valid after zero singular values are removed, or after restricting to the surviving full-row-rank submatrix.
